\section{On Hermetic Sublimation}

In \textbf{Valentin Tomberg}'s discussion of \textit{The Empress}, he quotes \textbf{Josephin Peladan}'s description of magic as ``\,`the sublimation of man'... of human nature''. Thereafter, Tomberg offers a detailed discussion differentiating Sacred Magic, from ``personal and arbitrary magic'', the latter which explains became more or less the earmark of Renaissance ``ceremonial magic'', of which he mostly looks askance.

This is of interest to the Hermetist for a number of reasons.

While \textbf{Marsilio Ficino} has been seen as ``reviving'' Hermetic magic in the Renaissance, it has been pointed out that this is a distortion stemming from Frances Yates' selective uses of Ficino's translation of \emph{Asclepius}, and not \emph{Pimander}. \textbf{Walter Haengraff} summarizes the situation:

\begin{quotex}
We are not dealing here with a straightforward case of a magical text influencing a later magical approach, but rather of an innovative interpretation of a non-magical text (i.e. The \emph{Corpus Hermetica}), resulting in a new perspective on how the attainment of a superior gnosis implies the acquisition of superhuman powers. Moreover the new perspective was not seen as magical either, but as something `better than magic'.
\end{quotex}


Thus, the Hermetica is not a ``magical'' text in the sense of what Tomberg calls ``personal and arbitrary magic'', but rather a scripture inviting to gnosis, the ``sublimation of man'', and as a by-product, certain ``powers'' might be arrived at---but not the other way around.

\textbf{Wouter Hanegraaff}, who is a scholar by occupation, has produced a critical edition of \textbf{Lodovico Lazzarelli}`s \emph{Crater Hermetis}---for the academic/scholarly background, readers are referred to that work, Hanegraaff's ``Better than Magic'' essay, and his related lecture series produced by the Ritman Library. In reference to Tomberg's meditation however, we only wish to consider certain of Lazzarelli's doctrines in ``Crater'', in direct relation with the idea of ``sublimation of man''. In his ``Crater'', Lazzerrelii (keeping in mind that he is considering himself to be a transmitter in the true Hermetic lineage) begins to introduce the central motif of his doctrine in 6.2. This could easily be considered ``divine procreation'', or ``sublimation'':

\begin{quotex}
Hermes says that God, having created all things in the beginning, exclaimed: ``Increase, grow up and multiply, all you seeds and works of my hands. And you, who have been given an inheritance of mind, recognize your origin and take heed of your immortal nature, and know that the love of body is the cause of death.' And that matches what Moses says in the book of Genesis. These words of Hermes contain the tree of life, by which we live, as well as the tree of knowledge of good and evil, which brings us death. And as you see, the main point of this precept is that we should know ourselves''.
\end{quotex}


It becomes quite interesting that the idea of ``liberation'' forms an element of Tomberg's meditation on Sacred Magic, the goal of which, similar to Lazzerrelli's above is ``... to follow an ideal and be truly oneself''. This sort of ``liberation'', which is a healing, Tomberg illustrates with example from the Gospel, where Christ produces miraculous cures. Perhaps a bit of synchronicity, or just broad overlap---but the second week of St. Ignatius' \textit{Spiritual Exercises}, focusing on the Via Illuminativa, the Kataphatic way, seeks to affirm everything that Peter denied; seeks to lead one into friendship with God, is intimately tied to the idea of ``liberation''---and likewise concentrates on the healings performed by Christ, such as making the hunchback woman upright (sublimated) on the Sabbath.

Lazzarrelli later in 21.4, continues his discussion concerning what the ``sublimation'' and ``self-knowledge'' engenders:

\begin{quotex}
Thus you will come to understand the excellence of your own essence, and will in no way disparage yourself, despise yourself, will not trample in your own supper, but you will rise up out of the body. Free from yourself and from all the things of the senses, ascending absolutely and purely, to fly up to that transcendent and most shining darkness where God dwells, to take your place among the number of Powers, and having been received among the Powers you shall enjoy God, and henceforth begetting a divine offspring, you will procreate for God and not for yourself. Because like is always produced by like''.
\end{quotex}


Consequently, as opposed to bells and whistles style ``ceremonial magic'', the ``Sacred Magic'', or Divine Gnosis to which the Hermetic Art invites us, might be more similar in practice to the Eastern idea of ``Bramacharya'' --- or, at least this appears to be what Lazzarelli is suggesting in his Hermetic transmission, as well as what we have read of Peladan above. In the applied Hermetism of Alchemy, just as on the Empresses' shield, is the Eagle, emblem of sublimation, distillation, which when ``flown'' after sufficient purgation, leads to the production of the Azoth, the ``Philosophical Mercury''.


\flrightit{Posted on 2014-03-13 by Frater M}

\begin{center}* * *\end{center}

\begin{footnotesize}
\begin{sffamily}
\texttt{nhm on 2014-03-13 at 22:55 said:}

The actual name of the scholar you are referring to is Wouter Hanegraaff. He is one of the editors of Dictionary of Gnosis and Western Esotericism, a book that I recommend wholeheartedly to all the readers of Gornahoor.

Thank you, corrections made. --- ed

\hfill

\texttt{JA on 2014-03-14 at 13:21 said:}

I am very much interested in reading actual practical Hermetic exercises based on the ideas presented by Tomberg and the other authors discussed here... ..

\end{sffamily}
\end{footnotesize}

